Under the completion compression calibration, the "radial deviation amount" corresponding to out-of-bound roots can be written as a completely explicit \(\mathrm{arcosh}\) closed form; this closed form compresses the horizontal offset of off-critical vertical lines from narrative description into a numerical invariant rigidly determined by the root positions of \(P_L\).\footnote{According to internal manuscript \cite{Internal2025ResolutionFoldingPhiPi} (2025), under the juxtaposition of "phase coordinate gate—finite window audit—orbit--Levi rigidity" without adding extra axioms, this deviation amount together with constant-layer gauge normalization forms an independently auditable closed chain. Archive timestamp: 2026-02-14 10:47:41 (Asia/Singapore).}

\begin{corollary}[Radial deviation closed form for single-scale completion: off-line distance for out-of-bound roots rigidly given by \texorpdfstring{$\mathrm{arcosh}$}{arcosh}]\label{cor:xi-single-scale-radial-deviation-arcosh}
Under the completion representation of Proposition \ref{prop:xi-rh-interval-criterion}, let
\[
\Xi_{\Delta,L}(s)=P_L\!\bigl(S_L(s)\bigr),\qquad S_L(s)=L^{s-\frac12}+L^{\frac12-s},
\]
where \(P_L\in\RR[S]\).
Take \(y:=L^{s-\frac12}\), then \(S_L(s)=y+y^{-1}\).
Let \(S_\star\in\RR\) be a root of \(P_L\) (counted with multiplicity), and let \(y_\star^{\pm}\) be the two solutions of the quadratic equation \(y+y^{-1}=S_\star\):
\[
y_\star^{\pm}=\frac{S_\star\pm\sqrt{S_\star^2-4}}{2}.
\]
Then by the multivaluedness of the complex logarithm, for each sign \(\pm\) and \(k\in\ZZ\) there is a zero tower
\[
s=\frac12+\frac{\log y_\star^{\pm}+2\pi i k}{\log L},
\]
and furthermore:
\begin{enumerate}
  \item If \(S_\star\in[-2,2]\), then \(|y_\star^\pm|=1\), thus all zeros induced by this root satisfy \(\Re(s)=\tfrac12\);
  \item If \(|S_\star|>2\), then the off-critical zeros induced by this root satisfy
  \[
  \bigl|\Re(s)-\tfrac12\bigr|
  =\frac{\mathrm{arcosh}\bigl(|S_\star|/2\bigr)}{\log L}.
  \]
\end{enumerate}
\end{corollary}

\begin{proof}
From \(S_L(s)=y+y^{-1}\) we obtain the quadratic equation \(y^2-S_\star y+1=0\), thus \(y=y_\star^\pm\).
Then from \(y=L^{s-\frac12}\) we get
\(
\log y=(s-\frac12)\log L-2\pi i k
\)
(\(k\in\ZZ\)), yielding the zero tower expression.
\par
When \(S_\star\in[-2,2]\) we can write \(S_\star=2\cos\theta\), then \(y_\star^\pm=e^{\pm i\theta}\), hence \(|y_\star^\pm|=1\).
When \(|S_\star|>2\), let \(|S_\star|=2\cosh\eta\) (\(\eta>0\), i.e., \(\eta=\mathrm{arcosh}(|S_\star|/2)\)), then the two roots satisfy \(|y_\star^\pm|=e^{\pm\eta}\).
From \(|y|=L^{\Re(s)-1/2}\) (proof of Proposition \ref{prop:xi-rh-interval-criterion}) we obtain
\(
|\Re(s)-\tfrac12|
=\eta/\log L
\), establishing the conclusion.
\end{proof}

